% ===== 原文 pp.145–155（PDF pp.152–162） =====
% [原文 p.145]
\chapter{有限维代数的表示论}

固定域 $K$，讨论 $K$ 上有限维代数中性质比较好的一类代数之上的模所成阿贝尔范畴与导出范畴的结构。

\section{有限维代数与其上的模}

固定域 $K$，向量空间、代数都在 $K$ 上考虑。

\begin{definition}{6.1}
域 $K$ 上的代数 $A$ 称为有限维 $K$ 代数（finite dimensional algebra），是指 $A$ 作为 $K$ 上的向量空间是有限维的。
\end{definition}

\begin{example}{6.2}
以 $K$ 为成分的 $n\times n$ 矩阵空间 $\mathrm{Mat}_n(K)$、下三角矩阵所成的部分 $K$ 代数（例 2.75）等都是有限维 $K$ 代数。多项式环 $K[x]$ 不是有限维 $K$ 代数。
\end{example}

设 $A$ 为有限维 $K$ 代数。$K$ 代数是环，因此对代数 $A$ 的（右、左）$A$ 模，可以按环 $A$ 上的 $A$ 模来定义。此时，每个有限生成（右）$A$ 模 $M$ 作为 $K$ 向量空间都是有限维的。事实上，取 $M$ 的生成元 $\xi_1,\dots,\xi_n$，则 $M$ 具有从自由 $A$ 模 $F := A^{\oplus n}$ 出发的满射 $\pi\colon F\to M$（命题 3.28）。$F$ 是有限维的，而 $\pi$ 是线性满射，故 $M$ 也是有限维的。进一步，对两个有限生成（右）$A$ 模 $M,N$ 之间的 $A$ 同态映射 $f\colon M\to N$，$\Ker(f),\Coker(f)$ 作为 $K$ 向量空间是有限维的，并且其作为 $K$ 向量空间的基可以取为右模的生成系，因此它们也是有限生成的。也就是说，$K$ 代数 $A$ 上的有限生成右 $A$ 模所成的加法范畴 $\modA$ 构成阿贝尔范畴。$\modA$ 的结构并不因为 $A$ 是有限维 $K$ 代数就变得简单，但存在若干类有限维 $K$ 代数 $A$，使得 $\modA$ 的结构可以被充分理解。特别地，下一章以后讨论作为箭图的（带关系式的）路代数而得到的有限维 $K$ 代数。已知这些代数给出了称为基本代数的这类代数的具体描述（定理 6.31）。下面先作一些相关术语的准备。

% [Tutor Note: 原书 p.145 此句作「次章以降では箙の（関係式付き）道代数として得られる有限次元 $K$ 代数について議論する」，直译「下一章以后」；但书中所指内容（箭图、路代数、带关系式的路代数）实为本第 6 章 6.2–6.4 节，作者此处「次章」当理解为「此后各节/此后」。照原印译出，特此说明。]

% [原文 p.146]
\begin{definition}{6.3}
$K$ 代数 $A$ 中满足 $e^2=e$ 的元 $e\in A$ 称为幂等元（idempotent）。特别地，当 $e$ 不能表示为与自身相异的幂等元 $e',e''$ 之和 $e=e'+e''$ 时，称 $e$ 为本原幂等元（primitive idempotent）。
\end{definition}

$0_A,1_A$ 是 $A$ 中的幂等元。

\begin{definition}{6.4}
在 $K$ 代数 $A$ 中，称 $1_A=e_1+\cdots+e_n$ 为 $1_A$ 的本原正交幂等元分解（primitive orthogonal idempotent decomposition），是指对每个 $i=1,\dots,n$，$e_i$ 都是本原幂等元，且 $i\neq j$ 时 $e_i\cdot e_j=0_A$。
\end{definition}

\begin{definition}{6.5}
称 $A$ 为 $K$ 上的基本代数（basic algebra），是指存在 $1_A$ 的本原正交幂等元分解 $1_A=e_1+\cdots+e_n$，且 $i\neq j$ 时右 $A$ 模 $e_iA$ 与 $e_jA$ 不同构。
\end{definition}

\begin{example}{6.6}
以域 $K$ 为成分的 $n\times n$ 矩阵环 $B=\mathrm{Mat}_n(K)$ 不是基本代数。事实上，取 $e_i\in B$ 为 $(i,i)$ 成分为 $1_K$、其他所有成分均为零的幂等元，则得到本原正交幂等元分解 $1_B=e_1+\cdots+e_n$，但对任意的 $i,j=1,\dots,n$，都得到右 $B$ 模的同构 $e_iB\simeq K^n\simeq e_iB$。此外，这个本原正交幂等元分解在 $B$ 的自同构下（除自同构外）是唯一的。

% [Tutor Note: 原书 p.146 例 6.6 印作「$e_iB\simeq K^n\simeq e_iB$」，第二处应为 $e_jB$（否则与「任意の $i,j$」及 $e_iB\simeq K^n$ 的论证不合）。作者正誤表（2014-05-20 版）未列此条，故正文照原印译出，存疑备考。]

另一方面，下三角矩阵所成的部分环 $A\subset B$ 构成基本代数。事实上，对本原正交幂等元分解 $1_A=1_B=e_1+\cdots+e_n$，对每个 $i$ 都得到右 $A$ 模的同构 $e_iA\simeq K^i$，且当 $i\neq j$ 时 $e_iA$ 与 $e_jA$ 作为 $K$ 向量空间不同构，因此作为右 $A$ 模也不可能同构。
\end{example}

给定基本 $K$ 代数 $A$ 与 $1_A$ 的本原正交幂等元分解 $1_A=e_1+\cdots+e_n$ 时，由于各 $e_i$ 是幂等元，$P(i):=e_iA$ 是投射模。对于 $1_A$ 的本原正交幂等元分解 $1_A=e_1+\cdots+e_n$，$A$ 是基本代数这一点意味着 $P(1),\dots,P(n)$ 恰好成为 $\modA$ 中不可分解投射对象的同构类的代表元列表；在这个意义上，范畴 $\modA$ 的结构论变得简明（参看 6.7 节）。

\section{箭图与路代数}

\begin{definition}{6.7}
箭图，或称 quiver（quiver）$\vec{\Delta}:=(\Delta_0,\Delta_1,s,t)$，是指集合 $\Delta_0$ 与 $\Delta_1$ 之组 $(\Delta_0,\Delta_1)$，并给定映射 $s\colon\Delta_1\to\Delta_0$，$t\colon\Delta_1\to\Delta_0$ 者。$\Delta_0,\Delta_1$ 的元分别称为箭图 $\vec{\Delta}$ 的顶点（vertex）、箭图 $\vec{\Delta}$ 的箭（arrow）。对箭 $\alpha\in\Delta_1$，称 $s(\alpha)$ 为 $\alpha$ 的起点（starting point, source），$t(\alpha)$ 为 $\alpha$ 的终点（ending point, target）。
\end{definition}

当 $\Delta_0,\Delta_1$ 都是有限集合时，称箭图 $\vec{\Delta}$ 是有限型的（finite type）。

由定义，可以说箭图 $\vec{\Delta}$ 就是由其顶点与射构成的有向图。

% [原文 p.147]
\begin{example}{6.8}
取箭图 $\vec{\Delta}=(\Delta_0,\Delta_1,s,t)$ 为 $\Delta_0=\{1,2\}$，$\Delta_1=\{\alpha\}$，$s(\alpha)=2$，$t(\alpha)=1$。这是有限型箭图，表示为如下有向图。
\[
\vec{\Delta} := \Bigl(\,
\begin{tikzpicture}[baseline=-0.5ex]
  \node[draw,circle,inner sep=1pt,label=below:{$1$}] (v1) at (0,0) {};
  \node[draw,circle,inner sep=1pt,label=below:{$2$}] (v2) at (1.8,0) {};
  \draw[-{Stealth}] (v2) -- node[above]{$\alpha$} (v1);
\end{tikzpicture}
\,\Bigr).
\]
\end{example}

\begin{definition}{6.9}
箭图 $\vec{\Delta}$ 中从顶点 $i\in\Delta_0$ 到顶点 $j\in\Delta_0$ 的长度 $r\ge 1$ 的路（path）$(i|\alpha_1\cdots\alpha_r|j)$，是指满足以下条件的 $r$ 个箭 $\alpha_1,\dots,\alpha_r\in\Delta_1$：
\begin{itemize}[leftmargin=2em,itemsep=0pt,topsep=2pt]
\item $s(\alpha_1)=i$，
\item $t(\alpha_k)=s(\alpha_{k+1})$（$k=1,\dots,r-1$），
\item $t(\alpha_r)=j$。
\end{itemize}
称 $s(w)=i$，$t(w)=j$ 分别为路 $w:=(i|\alpha_1\cdots\alpha_r|j)$ 的起点（starting point, source）、终点（ending point, target）。
\end{definition}

此外，对每个顶点 $i\in\Delta_0$，引入路 $(i|i)$，称其为从 $i$ 到 $i$ 的长度为零的路。

关于长度 $r\ge 1$ 的路的上述三个条件，无非是说：射 $\alpha_1$ 的起点为 $i$，射 $\alpha_k$ 的终点与 $\alpha_{k+1}$ 的起点一致，射 $\alpha_r$ 的终点为 $j$。由定义，长度为零的路只有 $\sharp(\Delta_0)$ 个。

\begin{definition}{6.10}
（参看 [2, chap.II.1]、[18, chap.1,5.1]、[16] 等。）箭图 $\vec{\Delta}$ 在 $K$ 上的路代数（path algebra）$K\vec{\Delta}$，是指以 $\vec{\Delta}$ 的路为基的 $K$ 向量空间，其积结构对基 $w:=(i|\alpha_1\cdots\alpha_r|j)$、$w':=(k|\beta_1\cdots\beta_{r'}|l)$ 由
\[
w\cdot w' := \delta_{jk}\,(i|\alpha_1\cdots\alpha_r\beta_1\cdots\beta_{r'}|l)
\]
确定。这里 $\delta_{jk}$ 仅在 $j=k$ 时为 $1$，其他情形为 $0$。
\end{definition}

也就是说，$K\vec{\Delta}$ 是具有与把两条路 $w,w'$ 连接起来相对应的积结构的代数，且规定当 $t(w)\neq s(w')$ 时 $w\cdot w'=0$。对每个 $i\in\Delta_0$，$e_i:=(i|i)$ 是 $A$ 中的本原幂等元；当 $\Delta_0$ 是有限集合时，$1_A:=\sum_{i\in\Delta_0}e_i$ 给出 $A$ 的幺元。也就是说，这给出了 $1_A$ 的本原正交幂等元分解。当 $\Delta_0$ 是无限集合时，关于 $e_i$ 的无限和不在 $A$ 中，故 $K\vec{\Delta}$ 是不具有幺元的 $K$ 代数。

给定箭图 $\vec{\Delta}$ 时，对每个 $\alpha\in\Delta_1$ 都存在长度 $1$ 的路 $(s(\alpha)|\alpha|t(\alpha))$。把对应的 $A:=K\vec{\Delta}$ 的基仍写作 $\alpha\in A$。同样地，把路 $w:=(i|\alpha_1\cdots\alpha_r|j)$ 与 $A$ 的基视为同一；利用 $A$ 中的积有 $w=\alpha_1\cdots\alpha_r$ 成立，因此关于 $A$ 的基，也采用这样的表示。

\begin{example}{6.11}
（下三角矩阵环）考虑如下有限型箭图 $\vec{\Delta}=(\Delta_0,\Delta_1,s,t)$。
\[
\vec{\Delta} := \Bigl(\,
\begin{tikzpicture}[baseline=-0.5ex]
  \node[draw,circle,inner sep=1pt,label=below:{$1$}] (v1) at (0,0) {};
  \node[draw,circle,inner sep=1pt,label=below:{$2$}] (v2) at (1.9,0) {};
  \node at (3.5,0) {$\cdots$};
  \node[draw,circle,inner sep=1pt,label=below:{$n$}] (vn) at (5.6,0) {};
  \draw[-{Stealth}] (v2) -- node[above]{$\alpha_{21}$} (v1);
  \draw[-{Stealth}] (3.0,0) -- node[above]{$\alpha_{32}$} (v2);
  \draw[-{Stealth}] (vn) -- node[above]{$\alpha_{n(n-1)}$} (4.1,0);
\end{tikzpicture}
\,\Bigr).
\]
\end{example}

% [原文 p.148]
也就是说，取 $\Delta_0:=\{1,2,\dots,n\}$，$\Delta_1:=\{\alpha_{(i+1)i}\mid i=1,\dots,n-1\}$，$s(\alpha_{(i+1)i})=i+1$，$t(\alpha_{(i+1)i})=i$。（取 $n=2$ 所得者给出例 6.8 的箭图。）此时，长度为零的路为
\[
(1|1),\qquad (2|2),\qquad \dots,\qquad (n|n)
\]
共 $n$ 个；长度为 $1$ 的路为
\[
(2|\alpha_{21}|1),\qquad (3|\alpha_{32}|2),\qquad \dots,\qquad (n|\alpha_{n(n-1)}|n-1)
\]
共 $n-1$ 个；长度为 $2$ 的路为
\[
(3|\alpha_{32}\alpha_{21}|1),\qquad (4|\alpha_{43}\alpha_{32}|2),\qquad \dots,\qquad (n|\alpha_{n(n-1)}\alpha_{(n-1)(n-2)}|n-2)
\]
共 $n-2$ 个；同样地反复进行，最后最长的（长度为 $n-1$ 的）路
\[
(n|\alpha_{n(n-1)}\cdots\alpha_{21}|1)
\]
有 $1$ 条。因此，路代数 $A:=K\vec{\Delta}$ 是以
\[
w_{ji}:=(j|\alpha_{j(j-1)}\cdots\alpha_{(i+1)i}|i),\qquad n\ge j\ge i\ge 1
\]
为基的 $K$ 向量空间
\[
A:=\Bigl\{\ \sum_{n\ge j\ge i\ge 1}a_{ji}w_{ji}\ \Bigm|\ a_{ji}\in K\ \Bigr\},
\]
且其积结构由 $w_{lk}\cdot w_{ji}=\delta_{kj}w_{li}$ 确定。这就是下三角矩阵环，即与下三角矩阵所成的 $K$ 向量空间
\[
\begin{pmatrix}
K & 0 & \cdots & 0\\
K & K & 0 & \cdots & 0\\
\vdots & \vdots & \ddots & & \vdots\\
K & K & \cdots & & K
\end{pmatrix}
:=
\Biggl\{
\begin{pmatrix}
a_{11} & 0 & \cdots & 0\\
a_{21} & a_{22} & 0 & \cdots & 0\\
\vdots & \vdots & \ddots & & \vdots\\
a_{n1} & a_{n2} & \cdots & & a_{nn}
\end{pmatrix}
\Biggr\}
\]
赋予矩阵的通常乘法所成的 $K$ 代数同构。

\begin{example}{6.12}
取有限型箭图 $\vec{\Delta}=(\Delta_0,\Delta_1,s,t)$ 为 $\Delta_0:=\{1,2,3,4\}$，$\Delta_1:=\{\alpha_{21},\alpha_{32},\alpha_{42}\}$，$s(\alpha_{ji})=j$，$t(\alpha_{ji})=i$。除 $4$ 条长度为零的路之外，$\vec{\Delta}$ 还具有路 $(2|\alpha_{21}|1)$、$(3|\alpha_{32}|2)$、$(4|\alpha_{42}|2)$、$(3|\alpha_{32}\alpha_{21}|1)$、$(4|\alpha_{43}\alpha_{21}|1)$。因此 $K\vec{\Delta}$ 可描述为对以下 $K$ 向量空间赋予由矩阵自然确定的乘法者。
\[
\vec{\Delta} =
\Biggl(
\begin{tikzpicture}[baseline=-0.5ex]
  \node[draw,circle,inner sep=1pt,label=left:{$1$}] (v1) at (0,0) {};
  \node[draw,circle,inner sep=1pt,label=right:{$2$}] (v2) at (2.2,0) {};
  \node[draw,circle,inner sep=1pt,label=right:{$3$}] (v3) at (3.1,1.5) {};
  \node[draw,circle,inner sep=1pt,label=right:{$4$}] (v4) at (3.1,-1.5) {};
  \draw[-{Stealth}] (v2) -- node[below]{$\alpha_{21}$} (v1);
  \draw[-{Stealth}] (v3) -- node[above left]{$\alpha_{32}$} (v2);
  \draw[-{Stealth}] (v4) -- node[below left]{$\alpha_{42}$} (v2);
\end{tikzpicture}
\Biggr),
\qquad
K\vec{\Delta} =
\begin{pmatrix}
K & 0 & 0 & 0\\
K & K & 0 & 0\\
K & K & K & 0\\
K & K & 0 & K
\end{pmatrix}.
\]

% [Tutor Note: 原书 p.148 例 6.12 的路列表中末项印作 $(4|\alpha_{43}\alpha_{21}|1)$，但 $\Delta_1=\{\alpha_{21},\alpha_{32},\alpha_{42}\}$ 中并无 $\alpha_{43}$，按上文 $(4|\alpha_{42}|2)$ 及右图，此处应为 $(4|\alpha_{42}\alpha_{21}|1)$。作者正誤表未列此条，正文照原印译出，存疑备考；右图与 $K\vec\Delta$ 矩阵（第 $4$ 行第 $2$ 列非零）均支持 $\alpha_{42}$ 读法。]
\end{example}

% [原文 p.149]
\begin{example}{6.13}
（自由 $K$ 代数（张量代数））取有限型箭图 $\vec{\Delta}_{(n)}=(\Delta_0,\Delta_1,s,t)$ 为 $\Delta_0:=\{o\}$，$\Delta_1:=\{\alpha_k\mid k=1,\dots,n\}$，$s(\alpha_k)=t(\alpha_k)=o$。
\[
\vec{\Delta}_{(n=1)}=\Bigl(\,
\begin{tikzpicture}[baseline=-0.5ex]
  \node[draw,circle,inner sep=1pt,label=below:{$o$}] (o) at (0,0) {};
  \draw[-{Stealth}] (o) to[out=135,in=45,looseness=5] node[above]{$\alpha_1$} (o);
\end{tikzpicture}
\,\Bigr),\qquad
\vec{\Delta}_{(n=2)}=\Bigl(\,
\begin{tikzpicture}[baseline=-0.5ex]
  \node[draw,circle,inner sep=1pt,label=below:{$o$}] (o) at (0,0) {};
  \draw[-{Stealth}] (o) to[out=115,in=65,looseness=6] node[above]{$\alpha_1$} (o);
  \draw[-{Stealth}] (o) to[out=245,in=295,looseness=6] node[below]{$\alpha_2$} (o);
\end{tikzpicture}
\,\Bigr),\qquad\cdots.
\]
此时，把箭 $\alpha_k$ 对不同的 $k$ 连接起来，可以作出任意长的路。长度为零的路只有 $(o|o)$ 一条。这个路代数 $K\vec{\Delta}$ 无非是以 $\alpha_k$，$k=1,\dots,n$ 为基的 $n$ 维 $K$ 向量空间 $V$ 的张量代数（例 2.76）。当 $n=1$ 时 $V$ 是 $1$ 维的，$K\vec{\Delta}\simeq K[x]$（单变量多项式环）。
\end{example}

\begin{definition}{6.14}
箭图 $\vec{\Delta}$ 中，从顶点 $v$ 到其自身的长度 $r\ge 1$ 的路称为有向环（oriented loop）。
\end{definition}

\begin{example}{6.15}
例 6.11、6.12 的箭图 $\vec{\Delta}$ 没有有向环。例 6.13 的箭图 $\vec{\Delta}_{(n)}$ 具有有向环。$(o|\alpha_1|o),\dots,(o|\alpha_n|o)$ 以及把它们连接起来的路都是有向环。
\end{example}

\begin{example}{6.16}
在下述箭图 $\vec{\Delta},\vec{\Delta}'$ 中，$\vec{\Delta}$ 没有有向环，而 $\vec{\Delta}'$ 具有 $(1|\alpha_{13}\alpha_{34}\alpha_{42}\alpha_{21}|1)$ 等有向环。
\[
\vec{\Delta} :=
\Biggl(
\begin{tikzpicture}[baseline=-0.5ex]
  \node[draw,circle,inner sep=1pt,label=left:{$1$}] (v1) at (0,2.2) {};
  \node[draw,circle,inner sep=1pt,label=right:{$2$}] (v2) at (2.2,2.2) {};
  \node[draw,circle,inner sep=1pt,label=left:{$3$}] (v3) at (0,0) {};
  \node[draw,circle,inner sep=1pt,label=right:{$4$}] (v4) at (2.2,0) {};
  \draw[-{Stealth}] (v2) -- node[above]{$\alpha_{21}$} (v1);
  \draw[-{Stealth}] (v3) -- node[left]{$\alpha_{31}$} (v1);
  \draw[-{Stealth}] (v4) -- node[right]{$\alpha_{42}$} (v2);
  \draw[-{Stealth}] (v4) -- node[below]{$\alpha_{43}$} (v3);
\end{tikzpicture}
\Biggr),\qquad
\vec{\Delta}' :=
\Biggl(
\begin{tikzpicture}[baseline=-0.5ex]
  \node[draw,circle,inner sep=1pt,label=left:{$1$}] (w1) at (0,2.2) {};
  \node[draw,circle,inner sep=1pt,label=right:{$2$}] (w2) at (2.2,2.2) {};
  \node[draw,circle,inner sep=1pt,label=left:{$3$}] (w3) at (0,0) {};
  \node[draw,circle,inner sep=1pt,label=right:{$4$}] (w4) at (2.2,0) {};
  \draw[-{Stealth}] (w2) -- node[above]{$\alpha_{21}$} (w1);
  \draw[-{Stealth}] (w1) -- node[left]{$\alpha_{13}$} (w3);
  \draw[-{Stealth}] (w4) -- node[right]{$\alpha_{42}$} (w2);
  \draw[-{Stealth}] (w3) -- node[below]{$\alpha_{34}$} (w4);
\end{tikzpicture}
\Biggr).
\]
\end{example}

\begin{definition}{6.17}
在箭图 $\vec{\Delta}=(\Delta_0,\Delta_1,s,t)$ 中，若集合 $\Delta_0$ 带有序 $\le$，且对两个顶点 $i,j\in\Delta_0$，当 $i\le j$ 时不存在满足 $s(w)=i$、$t(w)=j$ 的长度 $1$ 以上的路 $w$，则称 $\vec{\Delta}$ 为有序箭图（ordered quiver）\footnote{集合 $\Delta_0$ 上的序 $\le$ 是指满足以下条件的关系式 $\le$：对 $i,j,k\in\Delta_0$，\begin{itemize}[leftmargin=1.6em,itemsep=0pt,topsep=1pt]\item $i\le i$（自反律），\item 若 $i\le j$、$j\le i$ 则 $i=j$（反对称律），\item 若 $i\le j$、$j\le k$ 则 $i\le k$（传递律）。\end{itemize}}。
\end{definition}

以下把有限型箭图 $\vec{\Delta}$ 的顶点集合表示为 $\Delta_0:=\{1,2,\dots,n\}$。特别地，有限型有序箭图的情形，按序为 $1\le 2\le\cdots\le n$ 的方式给顶点标号。

\begin{example}{6.18}
例 6.11、6.12 的箭图 $\vec{\Delta}$ 是有序有限型箭图。
\end{example}

\begin{example}{6.19}
在例 6.16 中，箭图 $\vec{\Delta}'$ 不是有序箭图，而箭图 $\vec{\Delta}$ 是有序箭图。即使交换顶点 $\bullet_2,\bullet_3$ 的序，也仍是有序箭图。
\end{example}

注意：有限型箭图 $\vec{\Delta}$，当其没有有向环时，$\vec{\Delta}$ 就是有序箭图。

$K\vec{\Delta}$ 成为有限维 $K$ 代数（定义 6.1）的充分必要条件是 $\vec{\Delta}$ 为有限型有序箭图。事实上，若 $\Delta_0$ 是无限集合，则长度为零的路有无限多条。若 $\Delta_1$ 是无限集合，则至少有无限多条长度 $1$ 的路。另一方面，即使箭图 $\vec{\Delta}$ 是有限型的，只要存在有向环 $(i|\alpha_1\cdots\alpha_r|i)$，

% [原文 p.150]
那么对任意正整数，其 $n$ 重环，即长度为 $nr$ 的路
\[
(i|\ \overbrace{\underbrace{\alpha_1\cdots\alpha_r}_{r}\cdots\cdots\underbrace{\alpha_1\cdots\alpha_r}_{r}}^{n}\ |i)
\]
都存在。

由有限生成右 $K\vec{\Delta}$ 模全体所成的阿贝尔范畴 $\mathrm{mod}\text{-}K\vec{\Delta}$，由于射的空间成为 $K$ 向量空间，故构成 $K$ 上的阿贝尔范畴。同样地，$\mathrm{mod}\text{-}K\vec{\Delta}$ 的（上下）有界导出范畴 $\Db(\mathrm{mod}\text{-}K\vec{\Delta})$ 构成 $K$ 上的三角范畴。

% [Tutor Note: 原文此处只说了「$K\vec\Delta$ 有限维 ⟺ $\vec\Delta$ 有限型有序」的正反两面中的反面（无限顶点/无限箭 $\Rightarrow$ 无限维；有向环 $\Rightarrow$ 任意长路 $\Rightarrow$ 无限维），未写出正面「有限型且无有向环 $\Rightarrow$ 路长有界 $\Rightarrow$ 有限维」一句。正面由「无有向环的有限箭图，其路长 $\le\sharp\Delta_0-1$」立即得到；这是后面用「$n=\sharp\Delta_0$ 控制同调维数」的依据。]

\section{箭图的路代数的表示论}

对有限型箭图 $\vec{\Delta}$ 的路代数 $A:=K\vec{\Delta}$，讨论阿贝尔范畴 $\modA$ 的结构。特别地，存在称为箭图 $\vec{\Delta}$ 的表示所成阿贝尔范畴 $\mathrm{Rep}(\vec{\Delta})$ 者，并证明（作为加法范畴的）范畴等价
\[
\mathrm{Rep}(\vec{\Delta})\simeq\modA
\]
成立（[2, p72, Thm.1.6]、[5, p57, Thm.1.5] 等）。如此，$\mathrm{Rep}(\vec{\Delta})$ 给出了 $\modA$ 的等价而别种的描述；一般而言 $\mathrm{Rep}(\vec{\Delta})$ 更具体、更易理解。也就是说，用 $\mathrm{Rep}(\vec{\Delta})$ 的语言来考虑 $\modA$，使得通过线性代数的矩阵计算具体地进行不可分解 $A$ 模的分类等成为可能。

% [Tutor Note: 「直既約」本书指 indecomposable（不能写成两个非零对象之直和者），译为「不可分解」（与「单对象」simple 相区别）。$\mathrm{Rep}(\vec\Delta)\simeq\modA$ 的等价可用如下双向构造给出：一边是 $M\mapsto(\bigoplus_i Me_i,\ \text{各 }\alpha\text{ 的右乘})$，另一边是 $(V_i,B_\alpha)\mapsto\bigoplus_i V_i$ 并令 $\xi\cdot\alpha=\iota_{t(\alpha)}B_\alpha\pi_{s(\alpha)}(\xi)$，见下文 p.151–152 的 $\phi,\psi$。]

\begin{definition}{6.20}
箭图 $\vec{\Delta}$ 的表示（representation）$V:=(\bigoplus_{i\in\Delta_0}V_i,\bigoplus_{\alpha\in\Delta_1}B_\alpha)$，是指对箭图的每个顶点 $i\in\Delta_0$ 对应一个向量空间 $V_i$，对每个箭 $\alpha\in\Delta_1$ 对应一个线性映射
\[
B_\alpha\colon V_{s(\alpha)}\to V_{t(\alpha)}
\]
者。
\end{definition}

\begin{definition}{6.21}
设 $V:=(V,B)$ 与 $V':=(V',B')$ 为箭图 $\vec{\Delta}$ 的表示。若对每个 $i\in\Delta_0$ 所确定的线性映射 $f_i$ 之组
\[
f:=\{f_i\colon V_i\to V'_i\}_{i\in\Delta_0}
\]
对所有箭 $\alpha\in\Delta_1$ 满足交换性
\begin{equation}
B'_\alpha\circ f_{s(\alpha)}=f_{t(\alpha)}\circ B_\alpha \tag{6.1}
\end{equation}
则称 $f\colon V\to V'$ 为表示之间的射（morphism）。
\end{definition}

给定箭图 $\vec{\Delta}$ 时，$\vec{\Delta}$ 的表示与它们之间的射构成范畴。把它记作 $\mathrm{Rep}(\vec{\Delta})$。

% [原文 p.151]
构造函子 $\phi\colon\mathrm{Rep}(\vec{\Delta})\to\modA$。对箭图 $\vec{\Delta}$ 的表示 $V$，定义右 $A$ 模 $\phi(V)\in\modA$。$\phi(V)$ 作为 $K$ 向量空间定义为
\[
\phi(V):=\bigoplus_{i\in\Delta_0}V_i.
\]
把 $\phi(V)$ 到 $V_i$ 的射影（线性映射）记为 $\pi_i\colon\phi(V)\to V_i$，反之把自然的线性单射记为 $\iota_i\colon V_i\to\phi(V)$。此时 $\pi_i\circ\iota_i=\id_{V_i}$ 成立。$A$ 的右作用，只要决定 $A$ 的基对 $\phi(V)$ 的右作用便可确定。对 $\xi\in\phi(V)$、$e_k,\alpha\in A$，定义
\[
\xi\cdot\alpha:=\iota_{t(\alpha)}\circ B_\alpha\bigl(\pi_{s(\alpha)}(\xi)\bigr),\qquad
\xi\cdot e_k:=\iota_k\circ\pi_k(\xi).
\]
若这样定义 $e_k$ 的作用，则可知 $\xi\cdot 1_A=\xi$ 成立。另一方面，为使 $(\xi\cdot\alpha)\cdot\beta=\xi\cdot(\alpha\cdot\beta)$ 成立而定义右端，便得到
\[
\xi\cdot(\alpha\cdot\beta):=
\begin{cases}
\iota_{t(\beta)}\circ B_\beta\circ B_\alpha\bigl(\pi_{s(\alpha)}(\xi)\bigr), & t(\alpha)=s(\beta),\\[2pt]
0, & t(\alpha)\neq s(\beta).
\end{cases}
\]
注意本来当 $t(\alpha)\neq s(\beta)$ 时就有 $\alpha\cdot\beta=0$。同样地，对路 $w=\alpha_1\cdots\alpha_r$，定义
\[
\xi\cdot(\alpha_1\cdots\alpha_r):=\iota_{t(\alpha_r)}\circ B_{\alpha_r}\circ\cdots\circ B_{\alpha_1}\bigl(\pi_{s(\alpha_1)}(\xi)\bigr).
\]
以上定义了 $A$ 对 $\phi(V)$ 的右作用。也就是说，定义了 $\phi\colon\Ob(\mathrm{Rep}(\vec{\Delta}))\to\Ob(\modA)$。

% [Tutor Note: 底本 p.151 第 8 行印作 $\xi\cdot e_k:=\iota_{t(\alpha)}\circ\pi_{s(\alpha)}(\xi)$；据作者正誤表（2014-05-20 版）订正为 $\xi\cdot e_k:=\iota_k\circ\pi_k(\xi)$，此处按订正后的内容翻译。原印式中的 $t(\alpha),s(\alpha)$ 未指定 $\alpha$，且按此式不能得到 $\xi\cdot 1_A=\xi$（$e_k$ 是顶点 $k$ 对应的幂等元，自然应取 $\iota_k\circ\pi_k$）。可验证：$\sum_k\xi\cdot e_k=\sum_k\iota_k\pi_k(\xi)=\xi$ 即 $\xi\cdot 1_A=\xi$。]

对射 $f\colon V\to V'$，$V,V'\in\Ob(\mathrm{Rep}(\vec{\Delta}))$，令 $\phi(f)\in\Hom_{\modA}(\phi(V),\phi(V'))$ 为由 $f$ 自然诱导者即可。也就是说，对 $\xi\in\phi(V)$，由
\[
\phi(f)(\xi):=\sum_{i\in\Delta_0}\iota'_i\circ f_i\circ\pi_i(\xi)
\]
确定。它成为右 $A$ 模的映射，这一点由定义是显然的。以上定义了函子 $\phi\colon\mathrm{Rep}(\vec{\Delta})\to\modA$。注意由构造 $\phi$ 是加性的。

其次，构造加性函子 $\psi\colon\modA\to\mathrm{Rep}(\vec{\Delta})$。对有限生成右 $A$ 模 $M\in\modA$，考虑作为 $K$ 向量空间的直和分解
\[
M=\bigoplus_{k\in\Delta_0}M_k,\qquad M\cdot e_k=:M_k.
\]
把 $M_k$ 与前面构造中的 $V_k$ 视为同一，追溯那里构造的逆，从 $A$ 的右作用定义 $B_\alpha$，便得到箭图 $\vec{\Delta}$ 的表示 $V$。其次，对射 $f\colon M\to N$，$M,N\in\modA$，定义 $\psi(f)$。$f$ 对 $\xi\in M$、$a\in A$ 满足
\begin{equation}
f(\xi\cdot a)=f(\xi)\cdot a \tag{6.2}
\end{equation}

% [原文 p.152]
者。把 $f$ 表为 $f=\sum_{i,j\in\Delta_0}\iota'_j\circ f_{ij}\circ\pi_i$，其中 $f_{ij}:=\pi'_j\circ f\circ\iota_i$；则当 $a=e_i$、$\xi=\iota_j(v)$（$v\in V_j$）时，条件 (6.2) 为
\[
\delta_{ij}\sum_{k\in\Delta_0}f_{ik}(v)=\sum_{l\in\Delta_0}\delta_{il}f_{jl}(v)\ \bigl(=f_{ji}(v)\bigr),
\]
即知当 $i\neq j$ 时 $f_{ji}=0$。因此令 $f_{ii}=:f_i$，接着在条件 (6.2) 中取 $a=\alpha$、$\xi=\iota_{s(\alpha)}(v)$、$v\in V_{s(\alpha)}$ 来考虑，便得到
\[
f_{t(\alpha)}\circ B_\alpha(v)=B'_\alpha\circ f_{s(\alpha)}(v).
\]
这无非是 $\vec{\Delta}$ 的表示之间的射 $f=\{f_i\}_{i\in\Delta_0}$ 的条件式 (6.1)。由此得到函子 $\psi\colon\modA\to\mathrm{Rep}(\vec{\Delta})$。特别地，可知 $\phi,\psi$ 给出 $\mathrm{Rep}(\vec{\Delta})$ 与 $\modA$ 作为加法范畴的等价；由于 $\modA$ 是阿贝尔范畴，故 $\mathrm{Rep}(\vec{\Delta})$ 也是阿贝尔范畴。

\begin{remark}{6.22}
把箭图 $\vec{\Delta}$ 的表示及其间的射用矩阵表示，也许更易于理解。设 $\Delta_0=\{1,\dots,n\}$，对表示 $V:=\{\bigoplus_{i\in\Delta_0}V_i,\bigoplus_{\alpha\in\Delta_1}B_\alpha\}$，令
\[
B:=
\begin{pmatrix}
B_{11} & \cdots & B_{1n}\\
\vdots & \ddots & \vdots\\
B_{n1} & \cdots & B_{nn}
\end{pmatrix},
\qquad
B_{ji}:=\bigoplus_{\alpha\in\Delta_1,\ s(\alpha)=i,\ t(\alpha)=j}B_\alpha.
\]
特别地，当 $i=j$ 时令 $B_{ii}=0$。同样地，对表示 $V'$ 定义 $B'=\{B'_{ji}\}$。另一方面，把 $f\colon V\to V'$ 表为 $n\times n$ 对角矩阵：
\[
f:=
\begin{pmatrix}
f_1 & & \\
& \ddots & \\
& & f_n
\end{pmatrix}.
\]
此时，条件式 (6.1) 形式上看作
\[
B'\circ f=f\circ B.
\]
长度 $r$ 的路所对应的 $A$ 的元对 $M\in\modA$ 的作用的相容性，可通过注视 $(B')^r\circ f=f\circ(B)^n$ 的各成分来理解。
\end{remark}

% [Tutor Note: 原书 p.152 注意 6.22 末式印作 $(B')^r\circ f=f\circ(B)^n$，右端指数为 $n$；按上下文（同一长度 $r$ 的路所对应的 $B^r$ 与 $(B')^r$ 的相容性）应为 $(B)^r$。作者正誤表未列此条，正文照原印译出，存疑备考。]

在以下两个简单的箭图的例子中，试考虑对应的路代数 $A=K\vec{\Delta}$ 的有限生成右模所成阿贝尔范畴 $\modA$。

\begin{example}{6.23}
考虑例 6.8 的箭图
\[
\vec{\Delta}:=\Bigl(\,
\begin{tikzpicture}[baseline=-0.5ex]
  \node[draw,circle,inner sep=1pt,label=below:{$1$}] (v1) at (0,0) {};
  \node[draw,circle,inner sep=1pt,label=below:{$2$}] (v2) at (1.8,0) {};
  \draw[-{Stealth}] (v2) -- node[above]{$\alpha$} (v1);
\end{tikzpicture}
\,\Bigr)
\]
的路代数 $A:=K\vec{\Delta}$。已知 $A$ 与 $2\times 2$ 下三角矩阵所成的 $K$ 代数同构（例 6.11 中 $n=2$ 的情形）。

有限生成右 $A$ 模的分类，归结为对两个有限维向量空间 $V_1,V_2$ 的线性映射 $f\colon V_2\to V_1$ 的分类。对应于直和分解 $V_1\simeq\Coker(f)\oplus\Img(f)$、$V_2\simeq\Coim(f)\oplus\Ker(f)$ 的 $V:=\bigl(V_1\xleftarrow{\ f\ }V_2\bigr)\in\mathrm{Rep}(\vec{\Delta})$ 的直和分解为
\[
V=\bigl(0\xleftarrow{\ 0\ }\Ker(f)\bigr)\oplus
\bigl(\Img(f)\xleftarrow{\ \sim\ }\Coim(f)\bigr)\oplus
\bigl(\Coker(f)\xleftarrow{\ 0\ }0\bigr).
\]
\end{example}

% [原文 p.153]
得到。也就是说，令 $\dim_K(V_j)=n_j$（$j=1,2$），$\dim_K(\Ker(f))=k$，$\dim_K(\Img(f))=\dim_K(\Coim(f))=i$，则
\[
V\simeq\bigl(0\xleftarrow{\ 0\ }K\bigr)^{\oplus k}\oplus
\bigl(K\xleftarrow{\ 1\ }K\bigr)^{\oplus i}\oplus
\bigl(K\xleftarrow{\ 0\ }0\bigr)^{\oplus(n_2-i)}
\]
成为。作为结论，可知任意有限生成右 $A$ 模都由
\[
S(2):=\bigl(0\xleftarrow{\ 0\ }K\bigr),\qquad
P(1):=\bigl(K\xleftarrow{\ 0\ }0\bigr),\qquad
P(2):=\bigl(K\xleftarrow{\ 1\ }K\bigr)
\]
这 $3$ 种不可分解对象之直和得到。

% [Tutor Note: 底本 p.153 上式第三项的指数印作 $\oplus(n_2-i)$；按同一式第一、二项与 $V_2\simeq\Coim(f)\oplus\Ker(f)$、$V_1\simeq\Coker(f)\oplus\Img(f)$ 可知 $\dim\Coker(f)=n_1-i$，故该指数应为 $n_1-i$（验证：对第三项取维数，$\bigl(K\leftarrow 0\bigr)$ 的 $V_1$ 分量为 $K$，故总维数为 $n_2-i$，与 $\dim V_1=n_1$ 一般不符；而 $\dim V_1=i+(n_1-i)=n_1$ 才相符）。作者正誤表未列此条，正文照原印译出，存疑备考。]

$P(i)$，$i=1,2$ 是投射模。事实上 $P(i)$ 对应于作为右 $A$ 模由 $P(i)=e_i\cdot A$ 得到的对象。

由这些箭图表示的表示，可以直接读出 $S(2),P(1),P(2)$ 之间的射的空间。例如
\[
(1,0)\colon\quad
\begin{tikzcd}[row sep=2.2em, column sep=2.6em]
  P(1): & K \arrow[d,"1"'] & 0 \arrow[l,"0"'] \arrow[d,"0"] \\
  P(2): & K & K \arrow[l,"1"']
\end{tikzcd}
\qquad\qquad
(0,1)\colon\quad
\begin{tikzcd}[row sep=2.2em, column sep=2.6em]
  P(2): & K \arrow[d,"0"'] & K \arrow[l,"1"'] \arrow[d,"1"] \\
  S(2): & 0 & K \arrow[l,"0"']
\end{tikzcd}
\]
可知 $\Hom_{\modA}(P(1),P(2))\simeq K\simeq\Hom_{\modA}(P(2),S(2))$；由同样的计算有 $\Hom_{\modA}(P(2),P(1))=0=\Hom_{\modA}(S(2),P(2))$。

在 6.7.3 节中，将对该阿贝尔范畴 $\modA$、导出范畴 $\Db(\modA)$ 的结构作更详细的讨论。

\begin{example}{6.24}
（Kronecker 箭图（Kronecker quiver））考虑称为 Kronecker 箭图的以下箭图 $\vec{\Delta}$ 的路代数 $A$。
\[
\vec{\Delta}:=\Bigl(\,
\begin{tikzpicture}[baseline=-0.5ex]
  \node[draw,circle,inner sep=1pt,label=below:{$1$}] (v1) at (0,0) {};
  \node[draw,circle,inner sep=1pt,label=below:{$2$}] (v2) at (2.0,0) {};
  \draw[-{Stealth}] (v2) to[out=25,in=155] node[above]{$\alpha$} (v1);
  \draw[-{Stealth}] (v2) to[out=-25,in=-155] node[below]{$\beta$} (v1);
\end{tikzpicture}
\,\Bigr).
\]
此时，右 $A$ 模由箭图的表示 $V:=\bigl(V_1\xleftarrow[\ g\ ]{\ f\ }V_2\bigr)\in\mathrm{Rep}(\vec{\Delta})$ 所描述。令 $\dim_K(V_j)=n_j$（$j=1,2$），取 $V_1,V_2$ 的基构成同型 $V_j\simeq K^{n_j}$，并把对应于 $f,g$ 的表示矩阵仍记作 $f,g\in\mathrm{Mat}_{n_1\times n_2}(K)$。通过交换 $V_1,V_2$ 的基，$(f,g)$ 被某可逆矩阵 $P\in\mathrm{Mat}_{n_1\times n_1}(K)$、$Q\in\mathrm{Mat}_{n_2\times n_2}(K)$ 变换为 $P(f,g)Q:=(PfQ,PgQ)$。此时，称 $(f,g)$ 与 $(f',g')=P(f,g)Q$ 同型。也就是说，右 $A$ 模的同型类的分类归结为矩阵组 $(f,g)$ 的同型类的分类。关于这一分类的结果，依域 $K$ 的性质而有很大不同。当 $K=\C$（或 $K$ 是所谓代数闭域）时，通过用可逆矩阵 $P$ 与 $Q$ 作变换，例如可把 $f$ 取为

% [原文 p.154]
的形状。进一步，考虑在保持此 $f$ 不变的情况下把 $g$ 变换成尽可能简单的形状即可。这是也称为 Kronecker 问题的线性代数问题，可以完全分类 [16, p12, p20] [2]。
\end{example}

\section{带关系式的箭图的表示论}

任意有限维基本代数 $A$ 都被实现为在箭图 $\vec{\Delta}$ 的路代数中加入关系式者。在这种状况下，与关系式为平凡的 6.3 节的情形一样，可以考虑两个阿贝尔范畴。也就是说，首先有该箭图 $\vec{\Delta}$ 的带关系式 $\rho$ 的路代数 $A$ 上的有限生成右模所成的范畴 $\modA$；另一方面，对箭图与关系式之组 $(\vec{\Delta},\rho)$，定义了「$\vec{\Delta}$ 的带关系式 $\rho$ 的表示」所成的阿贝尔范畴。然后与 6.3 节一样，讨论这两个阿贝尔范畴的等价成立。

\subsection{箭图中的关系式}

\begin{definition}{6.25}
箭图的路代数 $K\vec{\Delta}$ 的关系式（relation），是指 $\vec{\Delta}$ 的长度为 $2$ 以上、且共有起点与终点的 $r$ 条路 $w_k$，$k=1,\dots,r$（所对应的 $K\vec{\Delta}$ 的基）的线性组合
\[
\rho:=\sum_{k=1}^{r}c_kw_k\in K\vec{\Delta}.
\]
称 $s(w_k),t(w_k)\in\Delta_0$ 分别为 $\rho$ 的起点（starting point, source）、终点（ending point, target），记作 $s(\rho),t(\rho)$。
\end{definition}

\begin{definition}{6.26}
对箭图 $\vec{\Delta}$ 的路代数 $A:=K\vec{\Delta}$，设给定有限个关系式 $\rho_1,\dots,\rho_d\in A$（$d\ge 1$）。由 $\rho:=\{\rho_1,\dots,\rho_d\}$ 生成的理想
\[
\cI_\rho:=A\cdot\rho_1\cdot A+\cdots+A\cdot\rho_d\cdot A
\]
所成的剩余环 $A/\cI_\rho$，称为箭图 $\vec{\Delta}$ 的带关系式的路代数（path algebra with relations）$(K\vec{\Delta};\rho_1,\dots,\rho_d)$。
\end{definition}

% [Tutor Note: 底本 p.154 第 $-6$ 行将上式印作 $\cI:=A\circ\rho_1\circ A+\cdots+A\circ\rho_d\circ A$（以合成号 $\circ$ 表积）；据作者正誤表（2014-05-20 版）订正为 $A\cdot\rho_1\cdot A+\cdots+A\cdot\rho_d\cdot A$，此处按订正后的内容翻译。此处 $\cdot$ 是 $K\vec\Delta$ 中的乘法（即路的连接），$\circ$ 会误导为态射的合成。]

\begin{example}{6.27}
考虑以下箭图 $\vec{\Delta}$。
\[
\vec{\Delta}:=
\begin{tikzpicture}[baseline=-0.5ex]
  \node[draw,circle,inner sep=1pt,label=below:{$1$}] (v1) at (0,0) {};
  \node[draw,circle,inner sep=1pt,label=below:{$2$}] (v2) at (1.9,0) {};
  \node[draw,circle,inner sep=1pt,label=below:{$3$}] (v3) at (3.8,0) {};
  \draw[-{Stealth}] (v2) -- node[above]{$\alpha_{21}$} (v1);
  \draw[-{Stealth}] (v3) -- node[above]{$\alpha_{32}$} (v2);
\end{tikzpicture}
\]
令 $\rho:=(3|\alpha_{32}\alpha_{21}|1)$，则带关系式的路代数 $(K\vec{\Delta},\rho)$ 是矩阵

% [原文 p.155]
\[
\begin{pmatrix}
K & 0 & 0\\
K & K & 0\\
0 & K & K
\end{pmatrix}
\]
所描述。注意 $(3,1)$ 成分为零。
\end{example}

\begin{example}{6.28}
考虑例 6.16 中的箭图 $\vec{\Delta}$。把 $K\vec{\Delta}$ 用矩阵表示，则如下。
\[
\vec{\Delta}:=
\Biggl(
\begin{tikzpicture}[baseline=-0.5ex]
  \node[draw,circle,inner sep=1pt,label=left:{$1$}] (v1) at (0,2.2) {};
  \node[draw,circle,inner sep=1pt,label=right:{$2$}] (v2) at (2.2,2.2) {};
  \node[draw,circle,inner sep=1pt,label=left:{$3$}] (v3) at (0,0) {};
  \node[draw,circle,inner sep=1pt,label=right:{$4$}] (v4) at (2.2,0) {};
  \draw[-{Stealth}] (v2) -- node[above]{$\alpha_{21}$} (v1);
  \draw[-{Stealth}] (v3) -- node[left]{$\alpha_{31}$} (v1);
  \draw[-{Stealth}] (v4) -- node[right]{$\alpha_{42}$} (v2);
  \draw[-{Stealth}] (v4) -- node[below]{$\alpha_{43}$} (v3);
\end{tikzpicture}
\Biggr),
\qquad
K\vec{\Delta}\simeq
\begin{pmatrix}
K & 0 & 0 & 0\\
K & K & 0 & 0\\
K & 0 & K & 0\\
K^2 & K & K & K
\end{pmatrix}.
\]
$(4,1)$ 成分的 $K^2$ 意味着起点为 $4$、终点为 $1$ 的路有 $2$ 条。在此加入关系式 $\rho:=(4|\alpha_{42}\alpha_{21}|1)-(4|\alpha_{43}\alpha_{31}|1)$，则在对应的带关系式路代数 $(K\vec{\Delta},\rho)$ 中成为 $(4|\alpha_{42}\alpha_{21}|1)=(4|\alpha_{43}\alpha_{31}|1)$。$(K\vec{\Delta},\rho)$ 作为矩阵为
\[
\begin{pmatrix}
K & 0 & 0 & 0\\
K & K & 0 & 0\\
K & 0 & K & 0\\
K & K & K & K
\end{pmatrix}
\]
所描述。$(4,1)$ 成分意味着对应于起点为 $4$、终点为 $1$ 的两条路的两个基被同一视，对应的商向量空间成为 $1$ 维。
\end{example}

\begin{example}{6.29}
考虑以下箭图 $\vec{\Delta}$。
\[
\vec{\Delta}:=
\begin{tikzpicture}[baseline=-0.5ex]
  \node[draw,circle,inner sep=1pt,label=below:{$1$}] (v1) at (0,0) {};
  \node[draw,circle,inner sep=1pt,label=below:{$2$}] (v2) at (1.8,0) {};
  \node[draw,circle,inner sep=1pt,label=below:{$3$}] (v3) at (3.6,0) {};
  \node[draw,circle,inner sep=1pt,label=below:{$4$}] (v4) at (5.4,0) {};
  \node[draw,circle,inner sep=1pt,label=below:{$5$}] (v5) at (7.2,0) {};
  \draw[-{Stealth}] (v2) -- node[above]{$\alpha_{21}$} (v1);
  \draw[-{Stealth}] (v3) -- node[above]{$\alpha_{32}$} (v2);
  \draw[-{Stealth}] (v4) -- node[above]{$\alpha_{43}$} (v3);
  \draw[-{Stealth}] (v5) -- node[above]{$\alpha_{54}$} (v4);
\end{tikzpicture}
\]
令 $\rho:=(4|\alpha_{43}\alpha_{32}|2)$，则带关系式路代数 $(K\vec{\Delta},\rho)$ 作为矩阵
\[
\begin{pmatrix}
K & 0 & 0 & 0 & 0\\
K & K & 0 & 0 & 0\\
K & K & K & 0 & 0\\
0 & 0 & K & K & 0\\
0 & 0 & K & K & K
\end{pmatrix}
\]
所描述。在 $(K\vec{\Delta},\rho)$ 中，除 $\rho\in K\vec{\Delta}$ 之外，由基 $\alpha_{54}\cdot\rho$、$\rho\cdot\alpha_{21}$、$\alpha_{54}\cdot\rho\cdot\alpha_{21}$ 张成的向量空间也变为平凡。

% [Tutor Note: 底本 p.155 第 $-5$ 行印作「基底 $\alpha_{54}\circ\rho,\ \rho\circ\alpha_{21},\ \alpha_{54}\circ\rho\circ\alpha_{21}$ で張られる」；据作者正誤表（2014-05-20 版）订正，将合成号 $\circ$ 改为乘点 $\cdot$（路代数中的乘法＝路的连接），此处按订正后的内容翻译。可验证：$\rho=(4|\alpha_{43}\alpha_{32}|2)$ 的起点为 $4$、终点为 $2$；$\alpha_{54}$ 的起点为 $5$ 终点为 $4$，$\alpha_{21}$ 的起点为 $2$ 终点为 $1$，故 $\alpha_{54}\cdot\rho$、$\rho\cdot\alpha_{21}$ 与 $\alpha_{54}\cdot\rho\cdot\alpha_{21}$ 依次是起点 $5$、终点 $4$／起点 $4$、终点 $2$／起点 $5$、终点 $1$ 的路，正是被关系式连带杀掉（化为 $0$）的那些基。]
\end{example}

\begin{example}{6.30}
（多项式环）例 6.13 中箭图 $\vec{\Delta}_{(n)}=:\vec{\Delta}$ 的路代数 $K\vec{\Delta}$ 是 $n$ 维向量空间 $V$ 的张量代数 $T(V)$。在其上加入以下关系式。
\[
\rho_{kl}:=\alpha_k\alpha_l-\alpha_l\alpha_k,\qquad k,l=1,\dots,n.
\]
\end{example}
